\documentclass[12pt]{article}

\usepackage[utf8]{inputenc}
\usepackage{amsmath, amssymb, amsthm}
\usepackage{geometry}
\geometry{a4paper, margin=2.5cm}

\title{Funcția Cumulativă de Repartiție și Teorema de Existență și Unicitate}
\author{}
\date{}

\theoremstyle{plain}
\newtheorem{theorem}{Teoremă}

\begin{document}

\maketitle

\section{Definiție}

Fie $(\Omega, \mathcal{F}, \mathbb{P})$ un spațiu de probabilitate.
Pentru o variabilă aleatoare reală $X : \Omega \to \mathbb{R}$, funcția
\textbf{cumulativă de repartiție} este:



\[
F_X(x) = \mathbb{P}(X \le x), \qquad x \in \mathbb{R}.
\]



\section{Proprietăți fundamentale}

Funcția de repartiție $F_X$ satisface:

\begin{enumerate}
    \item \textbf{Monotonie:}
    

\[
        x_1 \le x_2 \implies F_X(x_1) \le F_X(x_2).
    \]



    \item \textbf{Limite la infinit:}
    

\[
        \lim_{x \to -\infty} F_X(x) = 0, \qquad
        \lim_{x \to +\infty} F_X(x) = 1.
    \]



    \item \textbf{Continuitate la dreapta:}
    

\[
        \lim_{t \downarrow x} F_X(t) = F_X(x).
    \]



    \item \textbf{Salturile lui $F_X$:}
    

\[
        F_X(x) - F_X(x^-) = \mathbb{P}(X = x).
    \]


\end{enumerate}

\section{Teorema de existență și unicitate}

\begin{theorem}[Existență și Unicitate a măsurii asociate]
Fie $F : \mathbb{R} \to [0,1]$ o funcție care satisface:

\begin{enumerate}
    \item este monoton crescătoare;
    \item este continuă la dreapta;
    \item $\displaystyle \lim_{x \to -\infty} F(x) = 0$ și 
          $\displaystyle \lim_{x \to +\infty} F(x) = 1$.
\end{enumerate}

Atunci există o \textbf{unică} măsură de probabilitate 
$\mathbb{P}$ pe spațiul măsurabil $(\mathbb{R}, \mathcal{B}(\mathbb{R}))$ astfel încât:



\[
\mathbb{P}((-\infty, x]) = F(x), \qquad \forall x \in \mathbb{R}.
\]



Mai mult, această măsură este complet determinată de valorile lui $F$.
\end{theorem}

\section{Exemple}

Pentru o variabilă aleatoare continuă cu densitate $f$:


\[
F(x) = \int_{-\infty}^{x} f(t)\, dt.
\]



Pentru o variabilă aleatoare discretă cu valori $x_k$ și probabilități $p_k$:


\[
F(x) = \sum_{x_k \le x} p_k.
\]



\end{document}
